% --- Methods (cross-model advisory comparison) ---
% Adjust model_id strings and bundle/prompt versions to match your runs.

\paragraph{Cross-model advisory comparison.}
We compared three commercial large language models as \emph{offline} performance engineering advisors: \texttt{[Gemini model\_id]}, \texttt{[GPT-4 model\_id]}, and \texttt{[Claude model\_id]}.
Each model received an identical, frozen diagnostic bundle comprising representative application code for product search, user dashboard, recommendations, and checkout; \texttt{EXPLAIN (ANALYZE, BUFFERS)} plans for slow paths; a \texttt{pg\_stat\_statements} aggregate excerpt; slow-query log snippets; and a Locust workload summary (throughput and latency percentiles).
We used a single prompt template (version \texttt{[v]}, repository path \texttt{prompts/llm\_comparison\_task.md}) and required structured JSON output conforming to \texttt{llm\_comparison/schema/model\_output.schema.json}.
Sampling temperature was \texttt{[0.0--0.2]} and maximum output tokens \texttt{[N]} for all models; browsing or tool use was disabled unless uniformly enabled.
Recommendations were mapped to seven mutually exclusive primary categories: indexing, query rewrite, caching, batching, transaction handling, bottleneck diagnosis, and invalid/unsafe.
Expert evaluators first defined a gold recommendation set from the bundle \emph{without} viewing model outputs, then labeled each model-generated item as full, partial, or no match to gold, flagged invalid or unsafe content, and provided blinded Likert ratings for correctness, actionability, novelty/relevance, safety, and estimated implementation effort.

% --- Results ---
\paragraph{Recommendation overlap and categories.}
Table~\ref{tab:llm-category-counts} summarizes category counts.
Pairwise Jaccard similarity of deduplicated recommendation keys (Table~\ref{tab:llm-jaccard}) was \texttt{[low/moderate/high]}, indicating \texttt{[redundant vs complementary]} suggestions across models under identical evidence.

\paragraph{Quality and gold agreement.}
Table~\ref{tab:llm-human-agreement} reports precision against the expert gold set (full matches), partial and non-match rates, mean quality scores, and the rate of unsafe or invalid items.
Differences among models were \texttt{[small/moderate/large]} on these metrics for this single-bundle study.

\paragraph{Invalid and hallucination-like suggestions.}
Table~\ref{tab:llm-invalid-examples} lists representative failure modes (e.g., references to non-existent relations, dialect-incompatible syntax, or unsafe production changes without mitigation).
We report the fraction of items flagged as \texttt{invalid\_unsafe} in the structured output and during expert review.

% --- Discussion ---
\paragraph{Interpretation.}
Similarity on identical inputs supports a \emph{limited} claim: frontier models often surface overlapping \emph{classes} of remediation (indexes, batching, caching) for the same anti-patterns, but specific SQL and prioritization differ.
We do not equate advisory overlap with equivalence on unseen systems, prompts, or organizational constraints.

% --- Limitations ---
\paragraph{Limitations (cross-model study).}
The comparison uses one application snapshot and one prompt instantiation; statistical generality is therefore bounded.
We evaluated advisory text rather than autonomous runtime control.
Unless Table~\ref{tab:llm-implementation-yield} is populated by a controlled implementation arm with fixed effort per model, we do not attribute end-to-end latency improvements to a particular LLM.

% --- Minimal publishable sentence (Option B) ---
% Under identical diagnostic inputs and prompting, we report category distributions, pairwise Jaccard overlap on deduplicated items, precision against an expert-derived gold set, and invalid-suggestion rates for Gemini, GPT-4, and Claude as offline advisors; we do not claim model-agnostic optimization without additional bundles and optional controlled implementation.
